\documentclass[12pt]{article}
\usepackage{amsmath}
\usepackage{amsfonts}
\usepackage{amssymb}
\usepackage{physics}
\usepackage{geometry}
\usepackage{parskip}
\geometry{a4paper, margin=0.7in}


\title{02YMECH - Solution 8}
\date{Assigned for the week of Nov 10, 2025}
\author{}

\begin{document}
\maketitle

\section*{Solutions}

\begin{enumerate}

\item Let
\[
V(r)=-\frac{k}{r}, \qquad f(r)=-\frac{k}{r^2},
\]
and the orbit be an ellipse with semi–major axis \(a\) and eccentricity \(e\).
The polar equation is
\[
r(\theta)=\frac{a(1-e^2)}{1+e\cos\theta}.
\]

The angular momentum per unit mass is
\[
h=r^2\dot\theta=\sqrt{k a(1-e^2)}.
\]

The time average of any quantity \(A\) over one period \(T\) is
\[
\langle A\rangle=\frac{1}{T}\int_0^T A\,dt
=\frac{1}{T}\int_0^{2\pi} A\,\frac{r^2}{h}\,d\theta .
\]

\textbf{Average potential energy}

\[
\langle V\rangle
=-\frac{k}{Th}\int_0^{2\pi} r\,d\theta .
\]

Using the orbit equation,
\[
\int_0^{2\pi} r\,d\theta
=a(1-e^2)\int_0^{2\pi}\frac{d\theta}{1+e\cos\theta}
=2\pi a\sqrt{1-e^2}.
\]

Thus
\[
\langle V\rangle
=-\frac{2\pi k a\sqrt{1-e^2}}{Th}.
\]

The period is
\[
T=\frac{2\pi a^{3/2}}{\sqrt{k}},
\]
so
\[
\boxed{\langle V\rangle=-\frac{k}{a}}.
\]

\newpage

\textbf{Average kinetic energy}

The total energy of a Kepler orbit is
\[\ 
E=-\frac{k}{2a}.
\]

Since \(E=\langle T\rangle+\langle V\rangle\),
\[
\boxed{\langle T\rangle=\frac{k}{2a}}.
\]

\textbf{Comparison}

\[
\boxed{\langle T\rangle=-\tfrac12\langle V\rangle},
\]
which agrees with the virial relation, obtained here purely from explicit
time–average integrals.


\item Two masses \(m_1,m_2\) move in a circular orbit of radius \(a\) and period \(\tau\).
Using the reduced mass \(\mu\),

\[
\mu a\omega^2=\frac{G m_1 m_2}{a^2}, 
\qquad 
\omega=\frac{2\pi}{\tau},
\qquad
\frac{1}{\mu}=\frac{1}{m_1}+\frac{1}{m_2}.
\]

Hence
\[
a=\left(\frac{G m_1 m_2\,\tau^2}{4\pi^2\mu}\right)^{1/3}.
\]

After the motion is stopped, the reduced mass falls radially from rest at \(x=a\).
Energy conservation gives
\[
-\frac{G m_1 m_2}{a}
=\frac12\mu \dot x^2-\frac{G m_1 m_2}{x},
\]
so
\[
\dot x=\sqrt{\frac{2G m_1 m_2}{\mu}\left(\frac1x-\frac1a\right)}.
\]

The collision time is
\[
t=-\int_a^0
\frac{dx}{\sqrt{\tfrac{2G m_1 m_2}{\mu}\left(\frac1x-\frac1a\right)}}.
\]

Rewrite the integrand:
\[
t=-\sqrt{\frac{a\mu}{2G m_1 m_2}}
\int_a^0 \frac{\sqrt{x}}{\sqrt{a-x}}\,dx.
\]

Let \(x=y^2\) (\(dx=2y\,dy\)):
\[
\int_a^0 \frac{\sqrt{x}}{\sqrt{a-x}}\,dx
=2\int_{\sqrt a}^0 \frac{y^2}{\sqrt{a-y^2}}\,dy
=-\frac{\pi a}{2}.
\]

Therefore
\[
t=\frac{\pi a}{2}\sqrt{\frac{\mu}{2G m_1 m_2}}.
\]

Substitute \(a\) from above:
\[
t=\frac{\pi}{2\sqrt2}\sqrt{\frac{a^3}{G m_1 m_2/\mu}}
=\frac{\pi}{2\sqrt2}\frac{\tau}{2\pi}.
\]

\[
\boxed{\,t=\dfrac{\tau}{4\sqrt2}\,}
\]

Thus, after stopping the circular motion, the two particles collide after a time
\(\tau/(4\sqrt2)\).


\item Given
\[
\lambda(x)=k x^n,\qquad 0\le x\le L,\quad n>-1.
\]

\begin{enumerate}
\item[(a)] Total mass:
\[
M=\int_0^L kx^n\,dx=\frac{kL^{n+1}}{n+1}.
\]

\item[(b)] Center of mass:
\[
x_{\mathrm{CM}}=\frac{1}{M}\int_0^L x\,\lambda(x)\,dx
=\frac{n+1}{n+2}L.
\]

\item[(c)] Requiring \(x_{\mathrm{CM}}=\tfrac{3L}{4}\):
\[
\frac{n+1}{n+2}=\frac34 \;\Rightarrow\; n=2.
\]

For \(n=2\), the mass density increases quadratically with \(x\); the rod is strongly weighted toward the right end.
\end{enumerate}

\end{enumerate}


\end{document}
